\section{Operations on circuits on parties}
\label{sec:party_syntax_operations}

This section runs a circuit on parties beside untouched spectator factors, reads
a circuit without gates party by party, moves pair sources along a map of
parties, and reads a word on blocks of parties, on a subset of parties, and on
the parties of one cell. In each case the operator of the new circuit is the
operator of the old one tensored with an identity or relabelled along a
bijection of configurations, as in \cite[Section~5.1]{OpenAI2026PolynomialPEPS},
where an operation acts on some registers and leaves the others untouched.

Throughout, $P$ is a set of parties, a system is a family $X=(X_p)_{p\in P}$ of
finite basis sets, and, when $P$ is finite, $\mathcal H_X=\bigotimes_p\C^{X_p}$
has the configurations $X_P=\prod_pX_p$ as its basis. For a party circuit $c$
from $X$ to $Y$ whose gates have one outcome each and an interpretation
$\mathrm{ev}$ of its gates, $[c]$ denotes its operator from $\mathcal H_X$ to
$\mathcal H_Y$; for a word, the interpretation is that of pair sources. For a
region $D\subseteq P$ and a matrix $K$ from $\C^{X_D}$ to $\C^{Y_D}$, where $Y$
agrees with $X$ outside $D$, $K\otimes1$ is its placement between $X$ and $Y$
(Definition~\ref{def:sites_place_between}).

\subsection{Circuits beside spectator factors}
\label{sec:party_frame}

Let $W=(W_p)_{p\in P}$ be a system of finite basis sets, the spectators. The
system $W\times X$ has the basis $W_p\times X_p$ at $p$, and its configurations
on a region $D$ are paired with $W_D\times X_D$.

\begin{definition}[Circuits beside spectator factors]
    \label{def:party_gate_frame}
    \lean{QuantumCircuit.Party.GateFrame, QuantumCircuit.Party.GateFrame.IsCompatible,
        QuantumCircuit.Party.SourceGate.frame,
        QuantumCircuit.Party.SourceGate.toPairSource_frame,
        QuantumCircuit.Party.PairEffectGate.frame,
        QuantumCircuit.Party.PairEffectGate.toPairSource_frame,
        QuantumCircuit.Party.PartyCircuit.frame}
    \leanok
    \uses{def:party_circuit, def:party_pair_gates, def:party_weighted_gate,
        def:circuits_gate_family}
    A family of gates runs beside spectator factors when every gate $g$ on a cell
    $D$ from $X$ to $Y$ and every spectator system $W$ give a gate $g^W$ on $D$ from
    $W\times X$ to $W\times Y$, with the outcomes of $g^W$ in bijection with those
    of $g$. An interpretation $\mathrm{ev}$ is compatible with framing when, for
    every outcome $o$ of $g^W$ with corresponding outcome $o'$ of $g$,
    \begin{align}
        \mathrm{ev}(g^W,o)=1_{W_D}\otimes\mathrm{ev}(g,o'),
        \label{eq:party_frame_compatible}
    \end{align}
    read along the pairings $(W\times X)_D\cong W_D\times X_D$ and
    $(W\times Y)_D\cong W_D\times Y_D$. A pair source or a pair effect beside $W$
    keeps its endpoints, its fresh bases and its vector, the fresh factor placed in
    front of the spectator factor at each endpoint. The gates of monomials are
    framed as pair sources or pair effects, and a weighted gate on $D$ is framed by
    framing each of its monomials beside the spectators $(W_p)_{p\in D}$, with the
    same coefficients. The circuit $c^W$ beside $W$ is defined by induction on
    the circuit $c$: the empty circuit stays empty, a gate $g$ becomes $g^W$ on the
    same cell, a local operation $A$ at $p$ becomes $1_{W_p}\otimes A$ at $p$, local
    operations $A_p$ at every party become $1_{W_p}\otimes A_p$, and a composition
    becomes the composition of the framed circuits.
\end{definition}

\begin{theorem}[Framing keeps the syntax]
    \label{thm:party_frame_syntax}
    \lean{QuantumCircuit.Party.PartyCircuit.UsesOnly.frame,
        QuantumCircuit.Party.PartyCircuit.gateCells_frame,
        QuantumCircuit.Party.PartyCircuit.numGates_frame,
        QuantumCircuit.Party.PartyCircuit.gateCount_frame,
        QuantumCircuit.Party.PartyCircuit.participation_frame,
        QuantumCircuit.Party.PartyCircuit.participationCount_frame,
        QuantumCircuit.Party.PartyCircuit.frameLocationsEquiv,
        QuantumCircuit.Party.PartyCircuit.participants_frameLocationsEquiv,
        QuantumCircuit.Party.PartyCircuit.mapGates_frame,
        QuantumCircuit.Party.Word.sources_frame,
        QuantumCircuit.Party.Monomial.sources_frame,
        QuantumCircuit.Party.Monomial.effects_frame,
        QuantumCircuit.Party.Monomial.effectCount_frame,
        QuantumCircuit.Party.Monomial.sourceCount_frame,
        QuantumCircuit.Party.PartyCircuit.forallGate_isRight_frame}
    \leanok
    \uses{def:party_gate_frame, def:party_gate_counts, def:party_word,
        def:party_pair_gates, def:party_relabel, def:circuits_gate_morphism}
    Let $c$ be a party circuit and $W$ a spectator system. If $c$ uses only the
    parties of a set $R$, so does $c^W$. The circuit $c^W$ has the same list of
    gate cells, the same number of gates, the same gate count, and, at every party
    $p$, the same participation and the same number of gates containing $p$. Its
    gate occurrences are in bijection with those of $c$, corresponding occurrences
    having the same cell. If a morphism $\varphi$ of gate families satisfies
    $\varphi(g^W)=\varphi(g)^W$ for every gate $g$, then mapping $c^W$ along
    $\varphi$ gives the circuit beside $W$ of the image of $c$. A word beside $W$
    has the sources of the word. A monomial beside $W$ has the pair sources and
    the pair effects of the monomial, hence the same numbers of each, and a chain
    all of whose gates are pair effects stays such a chain.
\end{theorem}
\begin{proof}\leanok
    By induction on the circuit. Framing does not change the cell of a gate or the
    pair data of a source or an effect, and it adds no gate. For the parties
    outside $R$ at which the local operations $A_p$ are the identifications of
    equal bases $Y_p=X_p$, the matrices $1_{W_p}\otimes A_p$ are the
    identifications of $W_p\times Y_p=W_p\times X_p$. The list of gates of a
    composition is the concatenation of the two lists, so the lists of pair sources
    and pair effects of a monomial are those of the framed monomial.
\end{proof}

\begin{theorem}[Pair sources and pair effects beside spectators]
    \label{thm:party_frame_pair_compatible}
    \lean{QuantumCircuit.Party.SourceGate.regionMatrix_frame,
        QuantumCircuit.Party.PairEffectGate.regionMatrix_frame,
        QuantumCircuit.Party.SourceGate.isCompatible_frame,
        QuantumCircuit.Party.PairEffectGate.isCompatible_frame,
        QuantumCircuit.Party.chainEv_isCompatible_frame}
    \leanok
    \uses{def:party_gate_frame, def:party_pair_gates}
    The matrix of a pair source beside $W$ on its cell $D$ is $1_{W_D}\otimes M$,
    where $M$ is the matrix of the source on $D$, read along the pairings of
    configurations, and likewise for a pair effect. Hence the interpretations of
    pair sources, of pair effects and of the gates of monomials are compatible with
    framing.
\end{theorem}
\begin{proof}\leanok
    The cell $D$ consists of the two endpoints. The element of the matrix of the
    framed source at a pair of configurations vanishes unless the spectator parts
    agree at both endpoints, that is, on all of $D$, and it is then the element of
    $M$ at the remaining parts; this is the element of $1_{W_D}\otimes M$. The
    matrix of a pair effect is the adjoint of the matrix of the source with the same
    pair data, and the adjoint of $1_{W_D}\otimes M$, relabelled, is
    $1_{W_D}\otimes M^\dagger$, relabelled. A gate of a monomial is one or the
    other.
\end{proof}

\begin{theorem}[The operator of a circuit beside spectators]
    \label{thm:party_op_frame}
    \lean{QuantumCircuit.Party.PartyCircuit.op_frame,
        QuantumCircuit.Party.PartyCircuit.op_frame_unit,
        QuantumCircuit.Party.Word.op_frame}
    \leanok
    \uses{def:party_gate_frame, def:party_circuit, def:party_word}
    Let $P$ be finite, let the gates have one outcome each, and let $\mathrm{ev}$ be
    compatible with framing. For every party circuit $c$ from $X$ to $Y$ and every
    spectator system $W$,
    \begin{align}
        [c^W]=1_{W_P}\otimes[c],
        \label{eq:party_op_frame}
    \end{align}
    read along the pairings $(W\times X)_P\cong W_P\times X_P$ and
    $(W\times Y)_P\cong W_P\times Y_P$. In particular this holds for every word.
    If every $W_p$ has one element, then the element of $[c^W]$ at $(y,x)$ is the
    element of $[c]$ at the second components of $y$ and $x$.
\end{theorem}
\begin{proof}\leanok
    \uses{thm:party_frame_pair_compatible, lem:sites_place_between_products,
        thm:sites_product_span, def:circuits_single}
    By induction on the circuit. The empty circuit gives $1=1\otimes1$. For a gate
    $g$ on $D$, the operator of $g^W$ is the placement of
    $\mathrm{ev}(g^W,o)=1_{W_D}\otimes\mathrm{ev}(g,o')$
    by~(\ref{eq:party_frame_compatible}), the outcomes being unique, and placing
    $1_{W_D}\otimes K$ between $W\times X$ and $W\times Y$ gives
    $1_{W_P}\otimes(K\otimes1)$ by Lemma~\ref{lem:sites_place_between_products}. A
    local operation $1_{W_p}\otimes A$, read on the one-party region $\{p\}$, is
    $1_{W_{\{p\}}}\otimes A$ read there, and the same lemma applies. For local
    operations at every party, $\bigotimes_p(1_{W_p}\otimes A_p)
        =1_{W_P}\otimes\bigotimes_pA_p$ after regrouping the factors. For a
    composition, $(1\otimes[c'])(1\otimes[c])=1\otimes[c'][c]$. If every $W_p$ has
    one element, $1_{W_P}$ is the $1\times1$ identity.
\end{proof}

\begin{theorem}[Weighted gates beside spectators]
    \label{thm:party_frame_weighted}
    \lean{QuantumCircuit.Party.WeightedGate.isCompatible_frame}
    \leanok
    \uses{def:party_gate_frame, def:party_weighted_gate}
    The interpretation of weighted gates is compatible with framing.
\end{theorem}
\begin{proof}\leanok
    \uses{thm:party_op_frame, thm:party_frame_pair_compatible}
    The matrix of a weighted gate is $\sum_tc_t[M_t]$ over its terms, and that of the
    framed gate is $\sum_tc_t[M_t^W]$. By Theorem~\ref{thm:party_op_frame}, applied
    to the monomials with the compatible interpretation of their gates,
    $[M_t^W]=1\otimes[M_t]$, read along the pairings. By induction on the list of
    terms, the Kronecker product with $1$ and the relabelling are linear, which
    gives $1\otimes\sum_tc_t[M_t]$.
\end{proof}

\begin{theorem}[Framing keeps circuits allowed]
    \label{thm:party_frame_allowed}
    \lean{QuantumCircuit.Party.PartyCircuit.isAllowed_frame,
        QuantumCircuit.Party.Word.isAllowed_frame,
        QuantumCircuit.Party.PartyCircuit.IsLocalContractive.frame,
        QuantumCircuit.Party.ChainGate.isAllowed_frame,
        QuantumCircuit.Party.WeightedGate.isAllowed_frame}
    \leanok
    \uses{def:party_gate_frame, def:party_allowed, def:party_word,
        def:party_pair_gates, def:party_weighted_gate}
    Let $\mathrm{ev}$ be compatible with framing and let a condition on gates pass
    from every gate $g$ to $g^W$. If a party circuit is allowed for this condition,
    so is the circuit beside $W$. A circuit beside $W$ is locally contractive when
    the circuit is, and a word, an allowed gate of a monomial, and an allowed
    weighted gate stay allowed beside $W$.
\end{theorem}
\begin{proof}\leanok
    \uses{thm:detailed_balance_tensor_norm, lem:matrix_l2_contractions,
        thm:party_frame_pair_compatible}
    By induction on the circuit. A gate $g^W$ has by~(\ref{eq:party_frame_compatible})
    the matrix $1\otimes\mathrm{ev}(g,o')$, relabelled, whose norm is at most
    $\lVert1\otimes\mathrm{ev}(g,o')\rVert\le\lVert\mathrm{ev}(g,o')\rVert\le1$, and
    it satisfies the condition on gates by hypothesis. A local operation
    $1\otimes A$ has $\lVert1\otimes A\rVert\le\lVert A\rVert\le1$. For words and
    monomials the interpretations of pair sources and of the gates of monomials are
    compatible by Theorem~\ref{thm:party_frame_pair_compatible}, and the condition on
    their gates concerns the vectors, which framing keeps, so the first statement
    applies; every framed monomial of an allowed weighted gate is then allowed.
\end{proof}

\subsection{Circuits without gates}
\label{sec:party_local}

\begin{definition}[Circuits without gates and their party operations]
    \label{def:party_local_circuit}
    \lean{QuantumCircuit.Party.PartyCircuit.localOpFamily,
        QuantumCircuit.Party.PartyCircuit.localOpFamily_self,
        QuantumCircuit.Party.PartyCircuit.localOpFamily_of_ne,
        QuantumCircuit.Party.PartyCircuit.IsLocal,
        QuantumCircuit.Party.PartyCircuit.isLocal_iff_gates_eq_nil,
        QuantumCircuit.Party.Word.isLocal_iff_sources_eq_nil,
        QuantumCircuit.Party.PartyCircuit.partyOp,
        QuantumCircuit.Party.PartyCircuit.regionOp,
        QuantumCircuit.Party.PartyCircuit.parties}
    \leanok
    \uses{def:party_circuit, def:party_gate_counts, def:party_word,
        thm:sites_product_span}
    The party family of a local operation $A$ at $p$ from $X$ to $Y$, where
    $Y_s=X_s$ for $s\neq p$, is $A$ at $p$ and the matrix of the identification
    $X_s=Y_s$ at every $s\neq p$. A party circuit is local when it has no gate,
    that is, when its list of gates is empty; a word is local exactly when it has
    no pair source. The party operation $c_p$ of $c$ at $p$ is defined by induction
    on $c$: it is $1$ for the empty circuit, $0$ for a gate, the party family at $p$
    for a local operation, $A_p$ for local operations $(A_s)_s$ at every party, and
    $c'_pc_p$ for the composition of $c$ followed by $c'$. For a finite region $R$,
    the region operation is $c_R=\bigotimes_{p\in R}c_p$. When $P$ is finite, the
    parties of $c$ are those named by its operations: the cells of its gates, the
    party of each local operation, whatever its matrix, and, for local operations
    $(A_s)_s$ at every party, the parties $p$ at which $A_p$ is not the matrix of an
    identification of equal bases $Y_p=X_p$.
\end{definition}

\begin{theorem}[A circuit without gates acts party by party]
    \label{thm:party_local_op}
    \lean{QuantumCircuit.Party.PartyCircuit.op_localOp_eq_piKronecker,
        QuantumCircuit.Party.PartyCircuit.op_eq_piKronecker_partyOp,
        QuantumCircuit.Party.PartyCircuit.op_eq_reindex_kronecker_regionOp}
    \leanok
    \uses{def:party_local_circuit, def:sites_config}
    Let $P$ be finite and let the gates have one outcome each. The operator of one
    local operation is the product operator $\bigotimes_sm_s$ of its party family
    $(m_s)_s$. The operator of a local circuit $c$ is the product of its party
    operations,
    \begin{align}
        [c]=\bigotimes_{p\in P}c_p,
        \label{eq:party_local_op}
    \end{align}
    and, for every region $R$, in the coordinates of the cuts along $R$,
    $[c]=c_R\otimes\bigotimes_{p\notin R}c_p$. This is the collection of the local
    maps of a word without pair sources by party in the proof of
    \cite[Theorem~5.2]{OpenAI2026PolynomialPEPS}.
\end{theorem}
\begin{proof}\leanok
    \uses{def:party_circuit, def:sites_place_between, thm:sites_product_span,
        lem:sites_product_intertwine}
    The element of the placement of $A$ at $(y,x)$ is $A_{y_px_p}$ if $y$ and $x$
    agree outside $p$ through the identifications, and $0$ otherwise; the element
    of $\bigotimes_sm_s$ is $A_{y_px_p}\prod_{s\neq p}(m_s)_{y_sx_s}$, and the product
    over $s\neq p$ is $1$ exactly when $y$ and $x$ agree outside $p$. The identity
    (\ref{eq:party_local_op}) follows by induction on the circuit: the empty circuit
    gives $1=\bigotimes_p1$, local operations at every party have the product
    operator as their operator, and for a composition
    $(\bigotimes_pc'_p)(\bigotimes_pc_p)=\bigotimes_pc'_pc_p$. The last statement is
    (\ref{eq:party_local_op}) with the product split along the cut, by
    Lemma~\ref{lem:sites_product_intertwine}.
\end{proof}

\begin{theorem}[The party operations of an allowed circuit are contractions]
    \label{thm:party_local_contraction}
    \lean{QuantumCircuit.Party.PartyCircuit.norm_localOpFamily_le_one,
        QuantumCircuit.Party.PartyCircuit.norm_partyOp_le_one,
        QuantumCircuit.Party.PartyCircuit.norm_regionOp_le_one,
        QuantumCircuit.Party.PartyCircuit.exists_eq_piKronecker,
        QuantumCircuit.Party.Word.exists_piKronecker_of_sources_eq_nil}
    \leanok
    \uses{def:party_local_circuit, def:party_allowed, def:party_word}
    The party family of a contraction consists of contractions. If a party circuit
    is allowed for some condition on its gates, then every party operation $c_p$ and
    every region operation $c_R$ is a contraction. If $P$ is finite and $c$ is moreover
    local, then $[c]=\bigotimes_pB_p$ for contractions $B_p$. In particular an
    allowed word $w$ without pair sources satisfies $[w]=\bigotimes_pA_p$ for
    contractions $A_p$ \cite[Theorem~5.2]{OpenAI2026PolynomialPEPS}.
\end{theorem}
\begin{proof}\leanok
    \uses{lem:matrix_l2_contractions, thm:party_local_op}
    The matrix of an identification of bases is a contraction. By induction on
    the circuit, $c_p$ is $1$, $0$, a member of the party family of a contraction,
    the contraction $A_p$, or a product $c'_pc_p$ of two contractions. A product of
    contractions on the parties of $R$ is a contraction. For a local circuit take
    $B_p=c_p$ in (\ref{eq:party_local_op}); a word without pair sources is local.
\end{proof}

\begin{theorem}[The parties of a word]
    \label{thm:party_parties}
    \lean{QuantumCircuit.Party.Word.usesOnly_iff_parties_subset,
        QuantumCircuit.Party.Word.usesOnly_parties}
    \leanok
    \uses{def:party_local_circuit, def:party_word}
    Let $P$ be finite. A word uses only the parties of a set $R$ if and only if its
    parties lie in $R$. In particular a word uses only its own parties.
\end{theorem}
\begin{proof}\leanok
    By induction on the word. A source uses only $R$ when its cell lies in $R$, a
    local operation at $p$ when $p\in R$, local operations at every party when they
    are identifications of equal bases at every party outside $R$, and a
    composition when both parts do; the parties of a composition are the union of
    those of its parts.
\end{proof}

\subsection{Maps of source owners}
\label{sec:party_owner_map}

\begin{definition}[Pair sources along a map of parties]
    \label{def:party_owner_map}
    \lean{QuantumCircuit.Party.PairSource.mapOwner,
        QuantumCircuit.Party.SourceInventory.mapOwner,
        QuantumCircuit.Party.SourceInventory.mapOwner_nil,
        QuantumCircuit.Party.SourceInventory.mapOwner_cons,
        QuantumCircuit.Party.SourceInventory.mapOwner_append,
        QuantumCircuit.Party.affectedOwner,
        QuantumCircuit.Party.affectedOwner_eq_none,
        QuantumCircuit.Party.isSome_affectedOwner}
    \leanok
    \uses{def:party_pair_gates}
    Let $f\colon P\to Q$ be a map of parties. A pair source $s$ with endpoints $l$
    and $r$ is moved along $f$ to the pair source $f_*s$ with endpoints $f(l)$ and
    $f(r)$, the same fresh bases and the same vector, when $f(l)\neq f(r)$; when
    $f(l)=f(r)$ it becomes local and is dropped. A list $S$ of pair sources is moved
    to the list $f_*S$ of the moved sources that are kept, in their order, so that
    $f_*$ of the empty list is empty and $f_*(S\,S')=f_*S\,f_*S'$ for a concatenation.
    For a set $A\subseteq P$ of affected parties, the owner map $\omega_A$ from $P$
    to $A\sqcup\{\ast\}$ sends each affected party to itself and every other party to
    the one exterior party $\ast$; thus $\omega_A(p)=\ast$ if and only if
    $p\notin A$. This is the grouping of parties in the proof of
    \cite[Theorem~5.2]{OpenAI2026PolynomialPEPS}.
\end{definition}

\begin{theorem}[Moved lists of pair sources]
    \label{thm:party_owner_map}
    \lean{QuantumCircuit.Party.SourceInventory.mem_mapOwner,
        QuantumCircuit.Party.SourceInventory.IsNormalized.mapOwner,
        QuantumCircuit.Party.SourceInventory.mapOwner_eq_nil_iff,
        QuantumCircuit.Party.SourceInventory.map_partyPair_mapOwner,
        QuantumCircuit.Party.SourceInventory.mem_partyPair_mapOwner}
    \leanok
    \uses{def:party_owner_map}
    Let $f\colon P\to Q$ and let $S$ be a list of pair sources on $P$. A pair source
    $t$ occurs in $f_*S$ if and only if $t=f_*s$ for some $s$ in $S$ whose endpoints
    have distinct images. If every vector of $S$ is normalized, so is every vector
    of $f_*S$. The list $f_*S$ is empty if and only if the endpoints of every source
    of $S$ have the same image. The list of unordered endpoint pairs of $f_*S$ is
    obtained from that of $S$ by applying $f$ to each pair and dropping the pairs
    with equal points, in order and with multiplicity; hence an unordered pair
    $\{a,b\}$ occurs in it if and only if $a\neq b$ and $\{a,b\}$ is the image of an
    endpoint pair of $S$ \cite[Theorem~5.2]{OpenAI2026PolynomialPEPS}.
\end{theorem}
\begin{proof}\leanok
    The list $f_*S$ is obtained from $S$ by keeping the moved sources whose image
    endpoints differ, which gives the membership and emptiness statements; a moved
    source has the vector of the source. The list identity follows by induction on
    $S$, splitting on whether the endpoints of the first source have the same image.
\end{proof}

\begin{theorem}[Affected owners]
    \label{thm:party_affected_owner}
    \lean{QuantumCircuit.Party.affectedOwner_eq_iff,
        QuantumCircuit.Party.PairSource.affectedOwner_ne_iff,
        QuantumCircuit.Party.SourceInventory.mem_mapOwner_affectedOwner,
        QuantumCircuit.Party.SourceInventory.map_partyPair_affectedOwner,
        QuantumCircuit.Party.SourceInventory.mapOwner_affectedOwner_eq_nil_iff,
        QuantumCircuit.Party.affectedOwner_true_ne_none,
        QuantumCircuit.Party.uniqueBlockSystemNone}
    \leanok
    \uses{def:party_owner_map, def:sites_block}
    Let $A\subseteq P$. Two parties $p,q$ have the same owner, $\omega_A(p)=\omega_A(q)$,
    if and only if $p=q$ or both are unaffected. The endpoints of a pair source get
    distinct owners if and only if at least one of them is affected. Hence the
    sources occurring in $(\omega_A)_*S$ are exactly the moved sources $(\omega_A)_*s$
    of the sources $s$ of $S$ with an affected endpoint; its list of unordered
    endpoint pairs is the list of images under $\omega_A$ of the endpoint pairs of
    the sources of $S$ with an affected endpoint, in order and with multiplicity; and
    $(\omega_A)_*S$ is empty if and only if both endpoints of every source of $S$
    are unaffected. When every party is affected, no party is sent to $\ast$, and, for
    $P$ finite, the block of $\ast$ in a blocked system along $\omega_P$ has exactly
    one configuration \cite[Theorem~5.2]{OpenAI2026PolynomialPEPS}.
\end{theorem}
\begin{proof}\leanok
    \uses{thm:party_owner_map}
    Split on whether $p$ and $q$ are affected: two affected parties have the same
    owner exactly when they are equal, an affected party and an unaffected one never
    do, and two unaffected parties always do. The endpoints of a source are distinct,
    so they share an owner exactly when both are unaffected. The statements on
    $(\omega_A)_*S$ follow from Theorem~\ref{thm:party_owner_map}, and the list
    identity by induction on $S$. When every party is affected, the preimage of
    $\ast$ is empty and its configurations form the empty product.
\end{proof}

\subsection{Words along a map of parties}
\label{sec:party_coarsen}

Let $P$ be finite and $f\colon P\to Q$. The blocked system $X^f$ has the basis
$X_{f^{-1}(q)}$ at $q\in Q$ (Definition~\ref{def:sites_block}), and its
configurations are identified with those of $X$ by a bijection
$\theta_X\colon X^f_Q\to X_P$. A pair source $g$ with fresh bases $\alpha_g$ and
$\beta_g$ from $X$ to $Y$ identifies the configurations of its output with pairs,
$\gamma_g\colon Y_P\to(\alpha_g\times\beta_g)\times X_P$.

\begin{definition}[A word read on blocks of parties]
    \label{def:party_coarsen}
    \lean{QuantumCircuit.Party.freshBlockEquiv,
        QuantumCircuit.Party.freshBlockEquiv_apply_fst,
        QuantumCircuit.Party.freshBlockEquiv_apply_snd_self,
        QuantumCircuit.Party.freshBlockEquiv_apply_snd_of_ne,
        QuantumCircuit.Party.SourceGate.image_eq,
        QuantumCircuit.Party.SourceGate.coarse,
        QuantumCircuit.Party.SourceGate.toPairSource_coarse,
        QuantumCircuit.Party.Word.coarsen, QuantumCircuit.Party.Word.coarsen_comp,
        QuantumCircuit.Party.Word.coarsen_localOp}
    \leanok
    \uses{def:party_word, def:party_pair_gates, def:sites_block,
        lem:sites_subregion_place_between}
    Let $B$ be a region, $l\in B$, let $Y_v=X_v$ at the sites $v\in B$ other than $l$,
    and let $e\colon Y_l\to\alpha\times X_l$ be a bijection. Then a configuration $z$
    of $Y$ on $B$ is a pair of a fresh part $e(z_l)_1$ and the configuration of $X$ on
    $B$ equal to $e(z_l)_2$ at $l$ and to $z_v$ at $v\neq l$. A pair source $g$ on its
    cell $D$ has the image cell $f(D)=\{f(l),f(r)\}$. If $f(l)\neq f(r)$, the coarse
    source $g^f$ is the pair source on $f(D)$ from $X^f$ to $Y^f$ with endpoints
    $f(l)$ and $f(r)$, the fresh bases and the vector of $g$, its fresh factors located
    in the bases of the blocks by the preceding bijection. The word $w^f$ read on the
    blocks of $f$ is defined by induction on the word $w$:
    \begin{enumerate}
        \item the empty word stays empty, and a composition becomes the
              composition of the coarse words;
        \item a pair source $g$ with $f(l)\neq f(r)$ becomes $g^f$;
        \item a pair source $g$ with $f(l)=f(r)=b$ becomes one local operation at
              $b$, the matrix of $g$ on its cell extended to the block $f^{-1}(b)$ by
              the identity;
        \item a local operation $A$ at $p$ becomes a local operation at $f(p)$, the
              matrix $A$ read on the region $\{p\}$ and extended to $f^{-1}(f(p))$
              by the identity;
        \item local operations $(A_v)_v$ at every party become the local operations
              $\bigotimes_{v\in f^{-1}(b)}A_v$ at the blocks $b$.
    \end{enumerate}
    This is the reading on the affected and exterior systems in the proof of
    \cite[Theorem~5.2]{OpenAI2026PolynomialPEPS}, where sources inside one block
    become local preparations.
\end{definition}

\begin{theorem}[The coarse source on configurations]
    \label{thm:party_coarse_source}
    \lean{QuantumCircuit.Party.SourceGate.configEquiv_coarse}
    \leanok
    \uses{def:party_coarsen}
    Let $g$ be a pair source whose endpoints lie in distinct blocks. For every
    configuration $z$ of $Y^f$, with $(a,x)=\gamma_g(\theta_Y(z))$,
    \begin{align}
        \gamma_{g^f}(z)=\bigl(a,\theta_X^{-1}(x)\bigr).
        \label{eq:party_coarse_source}
    \end{align}
\end{theorem}
\begin{proof}\leanok
    The fresh parts agree by construction. Compare the old parts at a block $b$ and a
    site $v$ of $f^{-1}(b)$. If $b$ is the block of an endpoint and $v$ is that
    endpoint, both are the old part at the endpoint; if $v$ is another site of that
    block, it lies outside the cell, since the other endpoint is in another block, and
    both are $z$ at $v$. If $b$ is neither block, $v$ lies outside the cell and both
    are again $z$ at $v$.
\end{proof}

\begin{theorem}[The operator of a word read on blocks]
    \label{thm:party_op_coarsen}
    \lean{QuantumCircuit.Party.Word.op_coarsen, QuantumCircuit.Party.Word.sum_op_coarsen}
    \leanok
    \uses{def:party_coarsen}
    Let $Q$ be finite. For every word $w$ from $X$ to $Y$ and all configurations $z$
    of $Y^f$ and $z'$ of $X^f$,
    \begin{align}
        [w^f]_{z,z'}=[w]_{\theta_Y(z),\theta_X(z')}.
        \label{eq:party_op_coarsen}
    \end{align}
    For finitely many words $w_\xi$ and coefficients $c_\xi\in\C$, the combination
    $\sum_\xi c_\xi[w_\xi^f]$ is likewise $\sum_\xi c_\xi[w_\xi]$ read along
    $\theta_Y$ and $\theta_X$ \cite[Theorem~5.2]{OpenAI2026PolynomialPEPS}.
\end{theorem}
\begin{proof}\leanok
    \uses{thm:party_coarse_source, lem:sites_place_between_blocks,
        lem:sites_subregion_place_between, thm:sites_block_prod, def:party_word}
    By induction on the word. A pair source inside the block $b$ becomes the placement
    at $b$ of its extended matrix; by Lemma~\ref{lem:sites_place_between_blocks} this
    is, on the configurations of the parties, the extended matrix placed on $f^{-1}(b)$,
    which is the matrix of the source placed on its cell by
    Lemma~\ref{lem:sites_subregion_place_between}. A local operation is treated in the
    same way. The operator of a pair source joining two blocks is the preparation of
    its vector with the output read through $\gamma_g$, and
    (\ref{eq:party_coarse_source}) compares $\gamma_{g^f}$ with $\gamma_g$. Local
    operations at every party give, by Theorem~\ref{thm:sites_block_prod}, the product
    over the blocks of the products over the parties of each block. Products of
    relabelled matrices are relabelled products. The statement on linear combinations
    follows from (\ref{eq:party_op_coarsen}) entrywise.
\end{proof}

\begin{theorem}[Sources and allowedness of a word read on blocks]
    \label{thm:party_coarsen_sources}
    \lean{QuantumCircuit.Party.Word.sources_coarsen,
        QuantumCircuit.Party.Word.isAllowed_coarsen}
    \leanok
    \uses{def:party_coarsen, def:party_owner_map, def:party_word}
    The sources of $w^f$ are the sources of $w$ moved along $f$, $f_*$ of the list of
    sources of $w$. If $w$ is allowed, so is $w^f$.
\end{theorem}
\begin{proof}\leanok
    \uses{lem:sites_subregion_place_between, def:party_pair_gates}
    By induction on the word. A pair source inside one block becomes a local operation
    and is dropped by $f_*$; one joining two blocks becomes $g^f$, which is $f_*g$. The
    lists of a composition are concatenated and $f_*$ preserves concatenation. The
    extension of a matrix to a block does not increase its norm, the matrix of a pair
    source has norm at most that of its vector, and a product of contractions is a
    contraction.
\end{proof}

\subsection{Words restricted to a set of parties}
\label{sec:party_restrict}

\begin{definition}[A word restricted to a set of parties]
    \label{def:party_restrict}
    \lean{QuantumCircuit.Party.SourceGate.restrict,
        QuantumCircuit.Party.SourceGate.regionMatrix_restrict,
        QuantumCircuit.Party.Word.restrict, QuantumCircuit.Party.Word.numGates_restrict}
    \leanok
    \uses{def:party_word, def:party_pair_gates, def:sites_subregion_place}
    Let $S\subseteq P$ be finite. A pair source on a cell $D\subseteq S$ is read on
    the parties of $S$ with the same endpoints, fresh bases and vector; its matrix is
    the matrix of the source, read along the identification of the configurations of
    $D$ as a region of $P$ and as a region of $S$. A word $w$ that uses only the
    parties of $S$ restricts to the word $w|_S$ on the parties of $S$, from
    $(X_s)_{s\in S}$ to $(Y_s)_{s\in S}$, defined by induction on $w$: each pair
    source is read on $S$, a local operation at $p\in S$ is the same local operation
    at $p$, local operations at every party become those at the parties of $S$, and
    the empty word and compositions are kept. The word $w|_S$ has as many gates as
    $w$.
\end{definition}

\begin{theorem}[Sources and allowedness of a restricted word]
    \label{thm:party_restrict_sources}
    \lean{QuantumCircuit.Party.Word.sources_restrict,
        QuantumCircuit.Party.Word.isAllowed_restrict}
    \leanok
    \uses{def:party_restrict, def:party_owner_map}
    Let $w$ use only the parties of $S$. The sources of $w|_S$, moved along the
    inclusion of $S$ in $P$, are the sources of $w$. If $w$ is allowed, so is $w|_S$.
\end{theorem}
\begin{proof}\leanok
    \uses{lem:matrix_l2_contractions}
    By induction on the word. The inclusion is injective, so it keeps every source and
    moves it back to its endpoints in $P$; $f_*$ preserves concatenation. The matrix of
    a restricted source is a relabelling of the matrix of the source, which does not
    increase its norm, and its vector is unchanged.
\end{proof}

\begin{theorem}[A word using a set of parties is placed there]
    \label{thm:party_op_restrict}
    \lean{QuantumCircuit.Party.Word.op_eq_placeBetween_restrict}
    \leanok
    \uses{def:party_restrict, def:sites_place_between}
    Let $P$ be finite and let a word $w$ from $X$ to $Y$ use only the parties of a
    finite set $S$. Then $Y$ agrees with $X$ outside $S$, and
    \begin{align}
        [w]=[w|_S]\otimes1,
        \label{eq:party_op_restrict}
    \end{align}
    the operator of the restricted word placed on $S$ between $X$ and $Y$: the
    registers outside $S$ are untouched \cite[Section~5.1]{OpenAI2026PolynomialPEPS}.
\end{theorem}
\begin{proof}\leanok
    \uses{def:party_word, lem:sites_place_between_products,
        lem:sites_subregion_place_between, thm:party_local_op,
        lem:sites_place_between_algebra}
    By induction on the word. The empty word gives $1=\bigotimes_p1$, and placing the
    product operator of identities on $S$ gives it back, by
    Lemma~\ref{lem:sites_place_between_products}. A pair source on $D\subseteq S$ is
    the placement of its matrix on $D$, and by
    Lemma~\ref{lem:sites_subregion_place_between} this is the placement on $S$ of its
    extension from $D$ to $S$, which is the operator of the restricted source. A local
    operation at $p\in S$ is the product operator of its party family by
    Theorem~\ref{thm:party_local_op}, and so is its restriction; placing the latter on
    $S$ gives the product operator equal to the family on $S$ and to the
    identifications outside $S$, which is the former. Local operations at every party
    are the identifications of equal bases outside $S$, and the same comparison
    applies. A composition follows from
    $([w'|_S]\otimes1)([w|_S]\otimes1)=[w'|_S][w|_S]\otimes1$.
\end{proof}

\subsection{Words on the parties of a cell}
\label{sec:party_on_cell}

\begin{definition}[A word run on the parties of a cell]
    \label{def:party_on_cell}
    \lean{QuantumCircuit.Party.cellEquiv,
        QuantumCircuit.Party.cellEquiv_symm_apply_of_notMem,
        QuantumCircuit.Party.Word.onCell}
    \leanok
    \uses{def:party_coarsen, def:party_gate_frame, def:party_relabel}
    Let $D\subseteq P$ be finite, let $\iota\colon D\to P$ be the inclusion, and let
    $Y$ agree with $X$ outside $D$. The block of a party $p$ along $\iota$ is $\{p\}$
    if $p\in D$ and empty otherwise, and the spectator factor of $p$ is $X_p$ if
    $p\notin D$ and a one-point set otherwise. The basis $Y_p$ is identified with the
    product of the spectator factor and the basis of $p$ in the blocked system
    $(Y_s)_{s\in D}^\iota$; at $p\notin D$ the inverse identification reads the
    spectator factor through $Y_p=X_p$. A word $w$ on the parties of $D$, from
    $(X_s)_{s\in D}$ to $(Y_s)_{s\in D}$, is run inside all parties as the word
    $\widehat w$ from $X$ to $Y$ that first relabels $X$ along these identifications,
    then runs the word $w^\iota$ read on the blocks of $\iota$ beside the spectator
    factors, and finally relabels back to $Y$.
\end{definition}

\begin{theorem}[The operator of a word run on its cell]
    \label{thm:party_op_on_cell}
    \lean{QuantumCircuit.Party.Word.op_onCell,
        QuantumCircuit.Party.Word.isAllowed_onCell,
        QuantumCircuit.Party.Word.sources_onCell}
    \leanok
    \uses{def:party_on_cell, def:sites_place_between, def:party_owner_map}
    Let $P$ be finite. For every word $w$ on the parties of $D$,
    \begin{align}
        [\widehat w]=[w]\otimes1,
        \label{eq:party_op_on_cell}
    \end{align}
    the operator of $w$ placed on $D$ between $X$ and $Y$: an operation on the parties
    of a cell acts as the identity on the other parties
    \cite[Section~5.1]{OpenAI2026PolynomialPEPS}. The word $\widehat w$ is allowed
    when $w$ is, and its sources are those of $w$ moved along $\iota$.
\end{theorem}
\begin{proof}\leanok
    \uses{thm:party_op_frame, thm:party_op_coarsen, def:party_relabel,
        thm:party_frame_allowed, thm:party_coarsen_sources, thm:party_frame_syntax}
    The operator of $\widehat w$ is the product of the two relabellings with
    $[(w^\iota)^W]$, which is $1_{W_P}\otimes[w^\iota]$ by
    Theorem~\ref{thm:party_op_frame} and, by Theorem~\ref{thm:party_op_coarsen},
    $1_{W_P}\otimes[w]$ read on the configurations of $D$. Under the relabellings, the
    block part of a configuration $y$ of $Y$ is its restriction to $D$ and its
    spectator part is its restriction outside $D$, read in $X$. The element at $(y,x)$
    is therefore $[w]_{y|_D,x|_D}$ if $y$ and $x$ agree outside $D$ and $0$ otherwise,
    which is the element of $[w]\otimes1$. Relabellings are allowed and have no
    sources, framing keeps words allowed and keeps their sources by
    Theorems~\ref{thm:party_frame_allowed} and~\ref{thm:party_frame_syntax}, and
    reading on blocks does so up to moving the sources along $\iota$, by
    Theorem~\ref{thm:party_coarsen_sources}.
\end{proof}
